% =====================================================================
%  Commandes personnalisées
% =====================================================================

% ----- Placeholder figure (boîte TikZ avec libellé) -------------------
% Usage : \figplaceholder{Légende} ou \figplaceholder[hauteur=6cm]{Légende}
\NewDocumentCommand{\figplaceholder}{O{0.8\linewidth} O{6cm} m}{%
  \begin{tikzpicture}
    \draw[dashed,line width=0.6pt,fill=gray!8]
      (0,0) rectangle (#1,#2);
    \node[align=center,text width=0.9\dimexpr#1] at (#1/2,#2/2) {%
      \large\textbf{[FIGURE — à insérer]}\\[2pt]%
      \normalsize\itshape #3%
    };
  \end{tikzpicture}%
}

% ----- Notation cohérente VP/VN/FP/FN (directives Critères PDF) ------
\newcommand{\VP}{\mathrm{VP}}
\newcommand{\VN}{\mathrm{VN}}
\newcommand{\FP}{\mathrm{FP}}
\newcommand{\FN}{\mathrm{FN}}

% ----- Acronymes du projet ------------------------------------------
\newcommand{\FA}{FA}                       % Fibrillation auriculaire
\newcommand{\ECG}{ECG}
\newcommand{\CNN}{CNN}
\newcommand{\LSTM}{LSTM}
\newcommand{\AUROC}{AUROC}
\newcommand{\RR}{RR}

% ----- Encadré pour algorithmes (sans \fbox — directive 22 Sayadi) ----
% On utilise l'environnement `algorithm` qui est encadré par convention
% mais sans cadre dessiné autour du contenu.

% ----- Mot-clé visible (mise en valeur sans gras dans le texte) -------
\newcommand{\keyword}[1]{\emph{#1}}

% ----- Numérotation des équations par chapitre ----------------------
\numberwithin{equation}{chapter}
\numberwithin{figure}{chapter}
\numberwithin{table}{chapter}
